
\documentclass[a4paper,11pt]{article}
\usepackage{fontspec}
\setmainfont{Noto Sans}
\usepackage[french]{babel}
\usepackage{microtype}
\usepackage[top=9mm,bottom=8mm,left=11mm,right=11mm]{geometry}
\usepackage{xcolor}
\usepackage{tikz}
\usetikzlibrary{arrows.meta,calc,positioning}
\usepackage[most]{tcolorbox}
\usepackage{ulem}
\pagestyle{empty}

\definecolor{Blue}{HTML}{356FA8}
\definecolor{BlueDark}{HTML}{244D78}
\definecolor{BlueLight}{HTML}{EAF2F9}
\definecolor{BlueBorder}{HTML}{B9D3E3}
\definecolor{Violet}{HTML}{7A5AA6}
\definecolor{VioletDark}{HTML}{563A7A}
\definecolor{VioletLight}{HTML}{F2ECF8}
\definecolor{VioletBorder}{HTML}{D3C4E2}
\definecolor{Neutral}{HTML}{6C7378}
\definecolor{NeutralDark}{HTML}{42494D}
\definecolor{NeutralLight}{HTML}{F0F2F3}
\definecolor{NeutralBorder}{HTML}{C9CED1}
\definecolor{Ink}{HTML}{172A32}
\definecolor{Grey}{HTML}{64757D}
\definecolor{Red}{HTML}{D52F35}
\definecolor{RedLight}{HTML}{FCEBED}

\tcbset{
 enhanced,boxrule=.6pt,arc=2.8mm,
 left=3.7mm,right=3.7mm,top=1.75mm,bottom=1.75mm,
 boxsep=1mm,coltext=Ink
}
\newtcolorbox{bluebox}[1][]{colback=BlueLight,colframe=BlueBorder,#1}
\newtcolorbox{violetbox}[1][]{colback=VioletLight,colframe=VioletBorder,#1}
\newtcolorbox{neutralbox}[1][]{colback=NeutralLight,colframe=NeutralBorder,#1}
\newtcolorbox{warnbox}[1][]{colback=RedLight,colframe=Red,
 left=2.7mm,right=2.7mm,top=1.25mm,bottom=1.25mm,#1}

\begin{document}
\begin{tikzpicture}[remember picture,overlay]
\draw[draw=VioletBorder,line width=.9pt,rounded corners=1.7mm]
([xshift=1mm,yshift=-1mm]current page.north west)--([xshift=-1mm,yshift=-1mm]current page.north east)--
([xshift=-1mm,yshift=1mm]current page.south east)--([xshift=1mm,yshift=1mm]current page.south west)--cycle;
\fill[Violet]([xshift=1mm,yshift=-1mm]current page.north west) rectangle ([xshift=5.2mm,yshift=1mm]current page.south west);
\end{tikzpicture}
\vspace*{1mm}
\noindent\hspace*{6mm}{\color{VioletDark}\fontsize{25}{27}\selectfont\bfseries CONDICIONAL PRESENT}\hfill{\color{Grey}\large Le conditionnel présent}\par
\vspace{2.5mm}
\noindent\hspace*{6mm}\begin{minipage}{.91\textwidth}
\begin{violetbox}\textbf{\color{VioletDark}CONSTRUCCION}\\[1mm]\centering
{\Large\textbf{INFINITIU + TERMINASONS DEL CONDICIONAL}}\\[1mm]
\small \textit{cantar} $\rightarrow$ \textbf{cantariá, cantariás, cantariá, cantariam, cantaríatz, cantarían}
\end{violetbox}
\vspace{2.5mm}\begin{center}\begin{tikzpicture}[>=Latex]
\draw[Ink!65,line width=.8pt](0,0)--(15.1,0);
\node[below=2mm]at(2,0){\small PASSAT};\node[below=2mm]at(7.5,0){\small ARA};\node[below=2mm]at(13.5,0){\small FUTUR};
\draw[Violet,line width=2.5pt,->](7.8,0)--(12.8,0);
\draw[VioletDark,line width=1.1pt,->](7.5,1)--(7.5,.18);
\node[above=3.8mm]at(7.5,1){\bfseries\color{VioletDark}ARA};
\node[above=3mm,align=center,text width=6.2cm]at(10.3,0){\small\bfseries ipotèsi, desir, eventualitat, condicion};
\end{tikzpicture}\end{center}
\vspace{1.5mm}\begin{violetbox}\textbf{\color{VioletDark}A QUÉ SERVÍS ?}\\[1mm]
\begin{itemize}\setlength{\itemsep}{.4mm}
\item exprimir una \textbf{ipotèsi} o una condicion;
\item exprimir un desir, una possibilitat o una eventualitat;
\item formular una demanda o una suggestion de biais atenuat;
\item indicar una consequéncia ligada a una condicion.
\end{itemize}\end{violetbox}
\vspace{2mm}\begin{violetbox}\textbf{\color{VioletDark}EXEMPLES}\\[1mm]
\textbf{Se aviá mai de temps, viatjariá mai.}\\
\textbf{Voldriá visitar aquela vila.}\\
\textbf{Poiriás m'ajudar ?}\\
\textbf{Seriá interessant de veire aquò.}
\end{violetbox}
\vspace{2mm}
\begin{minipage}{.485\textwidth}\begin{violetbox}[height=3.15cm]
\textbf{\color{VioletDark}AMB « SE »}\\[1mm]\centering
\textbf{SE + IMPERFÈIT}\\
$\downarrow$\\
\textbf{CONDICIONAL}\\[1mm]\raggedright
\textit{Se aviá temps, vendriá.}
\end{violetbox}\end{minipage}\hfill
\begin{minipage}{.485\textwidth}\begin{warnbox}[height=3.15cm]
\textbf{\color{Red}ATENCION}\\[1mm]
Dins aquesta construccion, lo condicional apareis dins la proposicion de consequéncia, pas après \textbf{se}.
\end{warnbox}\end{minipage}
\vspace{2mm}\begin{violetbox}\textbf{\color{VioletDark}A RETÉNER}\\[1mm]
Lo condicional presenta una accion coma \textbf{ipotetica, possibla, desirabla o dependenta d'una condicion}.
\end{violetbox}
\vspace{1.5mm}\centering{\small\color{Grey}\textbf{Indicis frequents :} se, voldriá, poiriá, benlèu, a ta plaça, seriá...}
\vspace{1mm}{\scriptsize\color{Grey}Referéncia : occitan languedocian, graphia classica, vèrb'Òc / Lo Congrès.}
\end{minipage}\end{document}
